% GENERATED FILE. Edit canonical lessons, then run npm run generate:books.
% canonical-source-hash: fnv1a64:874ef2dbb26fbbec
% canonical-lessons: MW-C93-savan, MW-C93-bhadvo, MW-C93-asoj, MW-C93-kati, MW-C93-migsar

\chapter{Savan, Bhadvo, Asoj, Kati, Migsar}
\label{ch:mw-savan-bhadvo-asoj-kati-migsar}
\hlchaptermodality{93}

\begin{chapteropening}
\textbf{By the end of this chapter:} \emph{I can say Savan, Bhadvo, Asoj, Kati and Migsar.}

Complete the last lesson of chapter 93: I can say Savan, Bhadvo, Asoj, Kati and Migsar.
\end{chapteropening}

\section[\texorpdfstring{\mw{सावण} (sāvaṇ)}{सावण (sāvaṇ)}]{\mw{सावण} (sāvaṇ) — Savan (July-August)}

\label{lesson:MW-C93-savan}



\begin{quote}
Before the new one: say the Marwadi for Baisakh, then the Marwadi for Jeth.
\end{quote}

\subsection*{You'll want to know: \mw{सावण}}
\textbf{\mw{सावण}} — \emph{sāvaṇ} — \textquotedblleft{}Savan (July-August)\textquotedblright{}.

\begin{cousinweb}
\textbf{\mw{सावण}} (\emph{sāvaṇ}) is a tadbhava, a word inherited from Sanskrit and reshaped by sound change. It comes from \emph{śrāvaṇa}, \textquotedblleft{}the month of the Śravaṇa star\textquotedblright{}. CDIAL gives no Marwari form in this entry (CDIAL 12699); the match is through Hindi \emph{sāwan} and Gujarati \emph{sāvaṇ}.
\end{cousinweb}

\begin{grammarlens}[title={What you've built}]
One more word for this chapter.
\end{grammarlens}

\subsection*{Your turn}
\begin{itemize}
\raggedright
  \item \emph{Say it:} \emph{sāvaṇ}
  \item \emph{Say it:} \emph{sāvaṇ}, once more
  \item \emph{Say it:} say \emph{sāvaṇ}
  \item \emph{Recall:} say the Marwadi for Baisakh, then the Marwadi for Jeth, then say \emph{sāvaṇ} again
\end{itemize}

\begin{tcolorbox}[breakable,colback=teal!4,colframe=teal!35!black,title={Before you move on}]
What does \mw{सावण} mean? (\textquotedblleft{}Savan (July-August)\textquotedblright{}.) Say it once more.
\end{tcolorbox}
\section[\texorpdfstring{\mw{भादवो} (bhādvo)}{भादवो (bhādvo)}]{\mw{भादवो} (bhādvo) — Bhadvo (August-September)}

\label{lesson:MW-C93-bhadvo}



\begin{quote}
Before the new one: say the Marwadi for Jeth, then the Marwadi for Savan.
\end{quote}

\subsection*{You'll want to know: \mw{भादवो}}
\textbf{\mw{भादवो}} — \emph{bhādvo} — \textquotedblleft{}Bhadvo (August-September)\textquotedblright{}.

\begin{cousinweb}
\textbf{\mw{भादवो}} (\emph{bhādvo}) is a tadbhava, a word inherited from Sanskrit and reshaped by sound change. It comes from \emph{bhādrapada}, \textquotedblleft{}the month August-September\textquotedblright{}. CDIAL gives no Marwari form in this entry (CDIAL 9447); the match is through Hindi \emph{bhādõ} and Gujarati \emph{bhādarvo}.
\end{cousinweb}

\begin{grammarlens}[title={What you've built}]
One more word for this chapter.
\end{grammarlens}

\subsection*{Your turn}
\begin{itemize}
\raggedright
  \item \emph{Say it:} \emph{bhādvo}
  \item \emph{Say it:} \emph{bhādvo}, once more
  \item \emph{Say it:} say \emph{bhādvo}
  \item \emph{Recall:} say the Marwadi for Jeth, then the Marwadi for Savan, then say \emph{bhādvo} again
\end{itemize}

\begin{tcolorbox}[breakable,colback=teal!4,colframe=teal!35!black,title={Before you move on}]
What does \mw{भादवो} mean? (\textquotedblleft{}Bhadvo (August-September)\textquotedblright{}.) Say it once more.
\end{tcolorbox}
\section[\texorpdfstring{\mw{आसोज} (āsoj)}{आसोज (āsoj)}]{\mw{आसोज} (āsoj) — Asoj (September-October)}

\label{lesson:MW-C93-asoj}



\begin{quote}
Before the new one: say the Marwadi for Savan, then the Marwadi for Bhadvo.
\end{quote}

\subsection*{You'll want to know: \mw{आसोज}}
\textbf{\mw{आसोज}} — \emph{āsoj} — \textquotedblleft{}Asoj (September-October)\textquotedblright{}.

\begin{cousinweb}
\textbf{\mw{आसोज}} (\emph{āsoj}) is a tadbhava, a word inherited from Sanskrit and reshaped by sound change. It comes from \emph{aśvayuja}, \textquotedblleft{}the month Āśvina, September-October\textquotedblright{}. CDIAL gives no Marwari form in this entry (CDIAL 925); the match is through Hindi \emph{asoj}. The name is built on \emph{aśva}, \textquotedblleft{}horse\textquotedblright{}, and \emph{yuj}, \textquotedblleft{}yoke\textquotedblright{}, a star pair.
\end{cousinweb}

\begin{grammarlens}[title={What you've built}]
One more word for this chapter.
\end{grammarlens}

\subsection*{Your turn}
\begin{itemize}
\raggedright
  \item \emph{Say it:} \emph{āsoj}
  \item \emph{Say it:} \emph{āsoj}, once more
  \item \emph{Say it:} say \emph{āsoj}
  \item \emph{Recall:} say the Marwadi for Savan, then the Marwadi for Bhadvo, then say \emph{āsoj} again
\end{itemize}

\begin{tcolorbox}[breakable,colback=teal!4,colframe=teal!35!black,title={Before you move on}]
What does \mw{आसोज} mean? (\textquotedblleft{}Asoj (September-October)\textquotedblright{}.) Say it once more.
\end{tcolorbox}
\section[\texorpdfstring{\mw{काती} (kātī)}{काती (kātī)}]{\mw{काती} (kātī) — Kati (October-November)}

\label{lesson:MW-C93-kati}



\begin{quote}
Before the new one: say the Marwadi for Bhadvo, then the Marwadi for Asoj.
\end{quote}

\subsection*{You'll want to know: \mw{काती}}
\textbf{\mw{काती}} — \emph{kātī} — \textquotedblleft{}Kati (October-November)\textquotedblright{}.

\begin{cousinweb}
\textbf{\mw{काती}} (\emph{kātī}) is a tadbhava, a word inherited from Sanskrit and reshaped by sound change. It comes from \emph{kārttika}, \textquotedblleft{}the month October-November\textquotedblright{}. CDIAL gives no Marwari form in this entry (CDIAL 3070); the match is through Sindhi \emph{katī}. The month is named for the Kṛttikā stars, the Pleiades.
\end{cousinweb}

\begin{grammarlens}[title={What you've built}]
One more word for this chapter.
\end{grammarlens}

\subsection*{Your turn}
\begin{itemize}
\raggedright
  \item \emph{Say it:} \emph{kātī}
  \item \emph{Say it:} \emph{kātī}, once more
  \item \emph{Say it:} say \emph{kātī}
  \item \emph{Recall:} say the Marwadi for Bhadvo, then the Marwadi for Asoj, then say \emph{kātī} again
\end{itemize}

\begin{tcolorbox}[breakable,colback=teal!4,colframe=teal!35!black,title={Before you move on}]
What does \mw{काती} mean? (\textquotedblleft{}Kati (October-November)\textquotedblright{}.) Say it once more.
\end{tcolorbox}
\section[\texorpdfstring{\mw{मिगसर} (migsar)}{मिगसर (migsar)}]{\mw{मिगसर} (migsar) — Migsar (November-December)}

\label{lesson:MW-C93-migsar}



\begin{quote}
Before the new one: say the Marwadi for Asoj, then the Marwadi for Kati.
\end{quote}

\subsection*{You'll want to know: \mw{मिगसर}}
\textbf{\mw{मिगसर}} — \emph{migsar} — \textquotedblleft{}Migsar (November-December)\textquotedblright{}.

\begin{cousinweb}
\textbf{\mw{मिगसर}} (\emph{migsar}) is a tadbhava, a word inherited from Sanskrit and reshaped by sound change. It comes from \emph{mārgaśiras}, \textquotedblleft{}the month November-December\textquotedblright{}. CDIAL gives no Marwari form in this entry (CDIAL 10076); the match is through Hindi \emph{mãgsir} and Gujarati \emph{māgsar}.
\end{cousinweb}

\begin{grammarlens}[title={What you've built}]
One more word for this chapter.
\end{grammarlens}

\subsection*{Your turn}
\begin{itemize}
\raggedright
  \item \emph{Say it:} \emph{migsar}
  \item \emph{Say it:} \emph{migsar}, once more
  \item \emph{Say it:} say \emph{migsar}
  \item \emph{Recall:} say the Marwadi for Asoj, then the Marwadi for Kati, then say \emph{migsar} again
\end{itemize}

\begin{tcolorbox}[breakable,colback=teal!4,colframe=teal!35!black,title={Before you move on}]
What does \mw{मिगसर} mean? (\textquotedblleft{}Migsar (November-December)\textquotedblright{}.) Say it once more.
\end{tcolorbox}
